\section{Gr�ficos EPS}
\noindent
La tercera parte de la pr�ctica est� dedicada al manejo de gr�ficos. Para que
\LaTeX{} los pueda entender �stos deben estar en un formato adecuado.
La mejor calidad de impresi�n se obtiene con el formato \verb|.eps| (PostScript
Encapsulado, EPS). En principio se pueden utilizar otros formatos como
\verb|.png|, \verb|.tiff|, \verb|.jpeg|, aunque en realidad depende de la
plataforma y de la versi�n de \LaTeX{} que se puedan incluir o que la inclusi�n
sea m�s o menos sencilla.

\subsection{Gr�ficos EPS en su posici�n}
\noindent
Hay numerosas aplicaciones que generan ficheros .eps: Maple, Mathematica, MATLAB, Corel, etc.
Una vez que se tiene una imagen \verb|.eps| para incluirla en un fichero \LaTeX{} es
necesario utilizar macros que se encuentran en el paquete \verb|graphics|, e
incluir este paquete en el pre�mbulo del documento \verb|.tex| con la siguiente
instrucci�n: \begin{verbatim}\usepackage[driver]{graphicx}\end{verbatim} Donde
\verb|.driver| es el conversor de formato \verb|.dvi| formato \verb|.ps|. En
este caso el conversor es \verb|dvips|.\\
Esta informaci�n es necesaria porque los gr�icos los incluye en el fichero \verb|.dvi|
el driver de la impresora.\par
La orden para incluir un gr�fico ser�a:
\begin{verbatim}\includegraphics[clave=valor,...]{grafico}\end{verbatim}
Los pares clave-valor se emplean para modificar el ancho, la altura y el giro
del gr�fico. Las claves son:

\begin{itemize}
  \item \verb|height| (escalado al ancho)
  \item \verb|width| (escalado a la altura)
  \item \verb|angle| (giro)
\end{itemize}


�Cu�l ser�a el efecto de la siguiente orden?
\begin{verbatim}\includegraphics[width=4cm]{veritas.eps}\end{verbatim}
\includegraphics[width=4cm]{veritas.eps}

Veamos el efecto de la siguiente orden:
\begin{verbatim}\includegraphics[height=3cm]{veritas.eps}\end{verbatim}
\includegraphics[height=3cm]{veritas.eps}

Por �ltimo veamos el efecto de la siguiente orden:
\begin{verbatim}\includegraphics[height=5cm, angle=45]{veritas.eps}\end{verbatim}
\includegraphics[height=5cm, angle=45]{veritas.eps}

%A�adir otras dos aqu� con el formato especificado

%\includegraphics[height=5cm, angle=45]{veritas.eps}

%\includegraphics[height=5cm, angle=45]{veritas.eps}


\subsection{Gr�ficos EPS en elementos flotantes}
\noindent
En los ejemplos anteriores los gr�ficos han aparecido en la posici�n en la que
se han escrito las correspondientes �rdenes \verb|\includegraphics|.\par
Para no preocuparse de si el gr�fico va a entrar o no en la p�gina en el punto
en el que se hab�a escrito la orden de inclusi�n y que pueda quedar partido, lo
que se hace es dejar ``flotar'' el gr�fico hacia la siguiente p�gina.
La p�gina actual rellenar� el hueco que quede con texto. El elemento flotante
para la inserci�n de gr�ficos es \verb|figure|.\par

Veamos un ejemplo de uso de este elemento flotante:
\begin{verbatim}
\begin{figure}
 \includegraphics[width=7cm]{veritas.eps}
\end{figure}
\end{verbatim}
\begin{figure}
 \includegraphics[width=7cm]{veritas.eps}
\end{figure}

Veamos otro ejemplo de uso de este elemento flotante:
\begin{verbatim}
\begin{figure}
 \begin{center}
 \includegraphics[width=4cm, angle=90]{veritas.eps}
 \end{center}
\end{figure}
\end{verbatim}

\begin{figure}
 \begin{center}
 \includegraphics[width=4cm, angle=90]{veritas.eps}
 \end{center}
\end{figure}

Se pueden indicar preferencias de ubicaci�n:


\begin{itemize}
	\item ``h'' (``{\it here}''): aqu� o en la posici�n m�s pr�xima posible donde
	se encuentra declarada
	\item ``t'' (``{\it top}''): en la parte superior de una p�gina 
	\item ``b'' (``{\it bottom}''): en la parte inferior de una p�gina
	\item ``p'' (``{\it page}''): en una p�gina especial que s�lo contiene objetos
	flotantes (tablas y figuras)
\end{itemize}

Es razonable utilizar una referencia al elemento flotante. Con \verb|\label{}| se
proporciona una etiqueta a la figura y con \verb|\ref{}| se hace referencia a
la figura indicando la etiqueta anterior.\par
Un ejemplo de como etiquetar una figura para luego referenciarla es:

\begin{verbatim}
\begin{figure}
 \begin{center}
 \includegraphics[width=4cm, angle=180]{veritas.eps}
 \end{center}
 \caption{\label{figura1}Gr�fico de prueba}
\end{figure}
\end{verbatim}

\begin{figure}
 \begin{center}
 \includegraphics[width=4cm, angle=180]{veritas.eps}
 \end{center}
 \caption{\label{figura1}Gr�fico de prueba}
\end{figure}

Ahora se puede decir que la figura 1 presenta un gr�fico rotado 180 grados
(ejemplo de referencia a una figura).
